\documentclass[acmsmall,review,anonymous]{acmart}
\citestyle{acmnumeric}
\acmJournal{CSUR}
\acmDOI{}
\setcopyright{none}
\renewcommand\footnotetextcopyrightpermission[1]{}

\usepackage{amsmath,amssymb,amsfonts,mathtools}
\usepackage{booktabs,tabularx,multirow,array}
\usepackage{graphicx,xcolor,xspace}
\usepackage{enumitem}
\usepackage{url}

\newcolumntype{Y}{>{\raggedright\arraybackslash}X}
\newcommand{\RAG}{\textsc{RAG}\xspace}
\newcommand{\LLM}{\textsc{LLM}\xspace}
\newcommand{\QA}{\textsc{QA}\xspace}
\newcommand{\MH}{\textsc{MH}\xspace}
\newcommand{\CorpusSize}{232}
\newcommand{\AuditSize}{70}

\begin{document}

\title{Multi-Hop Retrieval and Reasoning: A Survey}

\author{Anonymous Author(s)}
\affiliation{%
  \institution{Anonymous Institution}
  \country{Anonymous}
}

\begin{abstract}
Multi-hop retrieval and reasoning systems solve information needs whose required evidence is distributed across sources and becomes identifiable only after intermediate results are resolved. The field now spans open-domain question answering, graph retrieval, retrieval-augmented generation, browser agents, deep research, long-form synthesis, and agents that improve through reflection, experience, or reusable skills. Existing surveys usually organize these areas by architecture, task, or data source. Such views are useful locally but make it difficult to compare static multi-hop pipelines with adaptive search agents, or to determine whether a method improves evidence discovery, subsequent search, evidence use, or future behavior.

We present a functional survey organized around four stages: \emph{initial retrieval}, \emph{follow-up retrieval}, \emph{evidence organization and reasoning}, and \emph{feedback learning}. The taxonomy treats deep search and deep research as compositions of these functions rather than separate top-level categories, and distinguishes feedback used within the current task from feedback that changes behavior across tasks. We formalize multi-hop systems as sequential evidence acquisition and transformation processes, provide crosswalks to retrieval, graph, agentic, memory, and skill-based architectures, and separate component claims from end-to-end answer accuracy. The working artifact contains a \CorpusSize-record literature corpus and a stratified \AuditSize-paper full-text audit selection, with source-verification status recorded separately. We conclude with an evaluation framework for retrieval depth, evidence coverage, context use, trajectory cost, grounding, and cross-task improvement.
\end{abstract}

\begin{CCSXML}
<ccs2012>
<concept>
<concept_id>10002951.10003317.10003347</concept_id>
<concept_desc>Information systems~Question answering</concept_desc>
<concept_significance>500</concept_significance>
</concept>
<concept>
<concept_id>10002951.10003317.10003338</concept_id>
<concept_desc>Information systems~Retrieval models and ranking</concept_desc>
<concept_significance>500</concept_significance>
</concept>
<concept>
<concept_id>10010147.10010178.10010179</concept_id>
<concept_desc>Computing methodologies~Natural language processing</concept_desc>
<concept_significance>500</concept_significance>
</concept>
</ccs2012>
\end{CCSXML}

\ccsdesc[500]{Information systems~Question answering}
\ccsdesc[500]{Information systems~Retrieval models and ranking}
\ccsdesc[500]{Computing methodologies~Natural language processing}

\keywords{multi-hop retrieval, multi-hop reasoning, retrieval-augmented generation, agentic search, deep research, evidence organization, experience learning, skill learning}

\maketitle

\section{Introduction}
\label{sec:introduction}

Many real-world questions cannot be answered by retrieving a single document. Their solutions require discovering intermediate entities, resolving hidden relations, and collecting information from sources that become relevant only after earlier evidence has been found. Consider the HotpotQA question~\cite{yang2018hotpotqa}:
\begin{quote}
\emph{What government position was held by the woman who portrayed Corliss Archer in the film \emph{Kiss and Tell}?}
\end{quote}
The original query does not directly name the target person. A system must first connect \emph{Kiss and Tell} to the character Corliss Archer, identify Shirley Temple as the actor who portrayed that character, and then search for the government position she later held. The information path is therefore approximately
\[
\textit{Kiss and Tell}
\rightarrow \textit{Corliss Archer}
\rightarrow \textit{Shirley Temple}
\rightarrow \textit{government position}
\rightarrow \textit{Chief of Protocol}.
\]
What makes the problem multi-hop is not merely that several documents are involved. The later information need is conditional on an earlier discovery: a search for Shirley Temple's government service becomes well specified only after the bridge entity has been resolved.

Conventional retrieval-augmented generation commonly follows a single-shot pattern,
\[
q \rightarrow \operatorname{Retrieve}(q) \rightarrow \operatorname{Generate}(q,E),
\]
which implicitly assumes that the original query is sufficient to retrieve the supporting evidence. This assumption is often violated in multi-hop settings. A more faithful abstraction is
\[
q_1=q, \qquad
E_t=\operatorname{Retrieve}(q_t), \qquad
q_{t+1}=\operatorname{Update}(q,E_{1:t}),
\]
followed by evidence integration and answer or report generation. The central challenge is therefore not only how to retrieve relevant information, but also how to decide what information is missing, what should be searched next, when the search should stop, and how evidence collected at different steps should be combined. Figure~\ref{fig:functional_process} shows the four functions and distinguishes adaptation within the current task from learning that persists across tasks.

\begin{figure*}[tbp]
\centering
\resizebox{0.98\textwidth}{!}{\begin{tikzpicture}[
  >=Latex,
  node distance=6mm,
  stage/.style={rounded corners=2pt, draw=#1!70!black, fill=#1!10, very thick,
                minimum height=13mm, text width=2.55cm, align=center, inner sep=4pt},
  endpoint/.style={rounded corners=2pt, draw=black!60, fill=black!4, thick,
                   minimum height=11mm, text width=1.65cm, align=center, inner sep=3pt},
  feedback/.style={rounded corners=2pt, draw=purple!65!black, fill=purple!8, thick,
                   text width=3.1cm, align=center, inner sep=4pt},
  arrow/.style={->, very thick, draw=black!65},
  loop/.style={->, thick, dashed, draw=#1!75!black},
  note/.style={font=\scriptsize, align=center, text=black!70}
]
\definecolor{initialblue}{RGB}{74,112,148}
\definecolor{followgreen}{RGB}{83,128,116}
\definecolor{integrateorange}{RGB}{168,119,71}
\definecolor{learnpurple}{RGB}{119,92,139}

\node[endpoint] (task) {Task\\$q$};
\node[stage=initialblue, right=of task] (initial) {\textbf{Initial Retrieval}\\Find an entry point\\$q\rightarrow E_1$};
\node[stage=followgreen, right=of initial] (follow) {\textbf{Follow-up Retrieval}\\Resolve the next need\\$(q,E_{1:t})\rightarrow E_{t+1}$};
\node[stage=integrateorange, right=of follow] (integrate) {\textbf{Evidence Integration}\\Select, preserve, and combine\\$E_{1:T}\rightarrow C$};
\node[endpoint, right=of integrate] (output) {Answer or\\cited report\\$y$};

\draw[arrow] (task) -- (initial);
\draw[arrow] (initial) -- node[above,note]{entry evidence} (follow);
\draw[arrow] (follow) -- node[above,note]{acquired sources} (integrate);
\draw[arrow] (integrate) -- (output);

\node[note, below=4mm of follow] (state) {Within-task state: entities, constraints, uncertainty,\\queries, observations, notes, and stopping signals};
\draw[loop=followgreen] (follow.south) to[out=-75,in=180] (state.west);
\draw[loop=followgreen] (state.east) to[out=0,in=-105] node[below right,note]{search again} (follow.south east);

\node[feedback, below=17mm of integrate] (learning) {\textbf{Feedback Learning}\\Diagnose trajectories; update memory, skills, retrievers, or policies\\$(\tau_i,f_i,M_i)\rightarrow M_{i+1},\pi_{i+1}$};
\draw[loop=learnpurple] (output.south) to[out=-90,in=25] node[right,note]{evaluation} (learning.east);
\draw[loop=learnpurple] (learning.west) to[out=165,in=-90] node[below left,note]{future tasks} (initial.south);
\draw[loop=learnpurple] (learning.north west) to[out=120,in=-90] (follow.south west);

\node[draw=black!25, rounded corners=2pt, fit=(task)(initial)(follow)(integrate)(output)(state),
      inner sep=4mm, label={[note]above:Current-task information acquisition}] {};
\node[draw=purple!35, rounded corners=2pt, fit=(learning),
      inner sep=3mm, label={[note]below:Cross-task adaptation}] {};
\end{tikzpicture}
}
\caption{Functional process of multi-hop retrieval and reasoning. Initial retrieval finds an entry point from the original task. Follow-up retrieval repeatedly resolves later information needs from the current state. Evidence integration selects, preserves, and combines the acquired sources into an answer or report. Within-task signals can redirect the current search, whereas feedback learning retains trajectory-level information to improve future retrieval, memory, skills, or policies.}
\Description{A left-to-right pipeline from task to initial retrieval, follow-up retrieval, evidence integration, and answer or report. A dashed within-task loop returns state information to follow-up retrieval. A lower feedback-learning box receives evaluation signals and updates future initial and follow-up retrieval.}
\label{fig:functional_process}
\end{figure*}

The field has evolved through several overlapping stages. Early multi-hop question answering benchmarks made cross-document dependencies explicit and often provided supporting-fact annotations~\cite{welbl2018qangaroo,yang2018hotpotqa,trivedi2022musique}. Multi-hop retrieval methods then learned to retrieve later passages from intermediate entities, paths, or previously retrieved passages~\cite{min2019golden,xiong2021mdr,khattab2021baleen}. Retrieval--reasoning systems subsequently interleaved search with generated reasoning or uncertainty signals~\cite{trivedi2023ircot,jiang2023flare,asai2024selfrag,jeong2024adaptiverag}. More recent agentic search and deep research systems extend the process to tool use, open-web navigation, longer horizons, source synthesis, and report generation~\cite{nakano2021webgpt,yao2023react,lala2023paperqa,jin2025searchr1}. These developments enlarge the operational setting, but they preserve the same core dependency: information found at one step determines what can be searched or inferred at the next.

Existing surveys provide valuable views of multi-hop question answering, retrieval-augmented generation, reasoning, and agentic systems~\cite{mavi2022mhqasurvey,gao2023ragsurvey,fan2024ragllmsurvey,singh2025agenticragsurvey}. Their organizing axes usually emphasize model architecture, retrieval mechanism, reasoning form, data source, or agent framework. These views do not always make it easy to compare systems that use different architectures to solve the same information-acquisition function. They also tend to treat evaluation primarily at the answer level, even when a method's main claim concerns query generation, hop completion, stopping, source use, trajectory efficiency, or improvement across tasks.

This survey adopts a functional perspective centered on the question:
\begin{quote}
\emph{How does a system find, follow, combine, and learn from information across multiple dependent hops?}
\end{quote}
We organize the literature into four categories. \emph{Initial retrieval} obtains entry evidence from the original query. \emph{Follow-up retrieval} determines and executes subsequent searches from intermediate results. \emph{Evidence integration} selects, preserves, organizes, and reasons over the collected information. \emph{Feedback learning} converts evaluations, trajectories, experience, or reusable procedures into improvements on future tasks. The first three functions primarily solve the current task; the fourth links completed tasks to future behavior. Deep search and deep research are consequently treated as system-level compositions of these functions rather than as independent top-level categories.

\paragraph{Scope.}
We focus on systems in which retrieval-grounded information acquisition is necessary for solving a multi-step task. The review covers textual, graph-structured, semi-structured, hybrid, and open-web evidence settings, together with long-form synthesis when it depends on multi-source acquisition. General-purpose agents, tool-use systems, memory architectures, and skill-learning methods are included only when they directly affect multi-hop retrieval, evidence integration, or cross-task improvement. The survey is therefore broader than a multi-hop RAG review, but it is not a general survey of autonomous agents.

\subsection{Contributions}
\label{subsec:contributions}

This survey makes four contributions.

\begin{itemize}[leftmargin=1.35em]
\item \textbf{A unified information-acquisition view.} We define multi-hop retrieval and reasoning as the process of acquiring, following, integrating, and learning from information across dependent steps. This definition separates genuine multi-hop dependence from tasks that merely contain multiple documents or multiple tool calls.

\item \textbf{A four-part functional taxonomy.} We organize methods into initial retrieval, follow-up retrieval, evidence integration, and feedback learning. The taxonomy provides a common language for comparing static multi-hop pipelines, graph-based systems, adaptive retrieval, agentic search, deep research, memory, and skill-based approaches without collapsing the survey into a generic agent taxonomy.

\item \textbf{A cross-generation review of methods and systems.} We connect early multi-hop QA and path retrieval with modern retrieval-augmented generation, graph retrieval, long-context evidence processing, browser-based search, deep research, and retrieval-grounded experience and skill learning. The companion artifact currently contains \CorpusSize{} literature records, with record-level source status exposed rather than treating all discovered papers as equally verified.

\item \textbf{A process-oriented evaluation framework.} We extend evaluation beyond final-answer accuracy to retrieval recall, hop-level success, query and action trajectories, evidence preservation and use, grounding, stopping behavior, cost, and cross-task improvement. A stratified \AuditSize-paper selection is used for the full-text diagnostic audit, with unreviewed fields left explicitly unresolved until source evidence is collected.
\end{itemize}

\paragraph{Organization.}
Section~\ref{sec:background} traces the evolution from multi-hop QA to deep research. Section~\ref{sec:overview} introduces the basic concepts, formalization, running examples, and taxonomy. Sections~\ref{sec:initial_retrieval}--\ref{sec:feedback_learning} review the four functional categories. Section~\ref{sec:evaluation} organizes evaluation by component and system claim. Section~\ref{sec:future} presents open problems, and Section~\ref{sec:conclusion} concludes the survey. The appendices document the review protocol, datasets, full taxonomy, and diagnostic audit.

\input{sections/02_methodology}
\section{Overview}
\label{sec:overview}

This section defines the unit of analysis used throughout the survey. The goal is to describe multi-hop systems at a level that is independent of a particular retriever, language model, graph representation, or agent framework.

\subsection{Basic Concepts}
\label{subsec:basic_concepts}

\paragraph{Information need.}
An information need is a statement of what must be resolved to advance the task. It may be expressed as a textual query, a graph traversal condition, a tool call, or an implicit retrieval decision.

\paragraph{Evidence unit.}
An evidence unit is the smallest retrievable or usable item in a system's environment, such as a passage, sentence, table row, knowledge-graph edge, image region, web page, or previously generated note. The appropriate granularity depends on the retrieval and integration mechanism.

\paragraph{Hop.}
A hop is a state transition in which acquired information changes the next retrieval or reasoning decision. Multiple retrieved documents do not imply multiple hops if they are obtained independently from the same unchanged query. Conversely, a single document can contain several dependent reasoning steps.

\paragraph{Bridge information.}
Bridge information resolves an intermediate entity, relation, constraint, or subproblem that makes later evidence identifiable. A bridge may be explicitly annotated in a benchmark or inferred from the system trajectory.

\paragraph{Trajectory.}
A trajectory records the sequence of states, queries, retrieval actions, observations, reasoning outputs, tool calls, stopping decisions, and final outputs produced for one task. Trajectories are central for evaluating follow-up retrieval and feedback learning.

\paragraph{Evidence integration.}
Evidence integration transforms the acquired evidence set into a reader- or generator-facing representation. It includes evidence selection, compression, ordering, structured context construction, cross-evidence reasoning, and long-form synthesis.

\paragraph{Feedback.}
Within-task feedback modifies a later action in the current trajectory, for example when low confidence triggers another retrieval. Cross-task feedback is retained after the current task and changes future behavior through reflection, memory, skill extraction, retriever updates, or policy learning.

\subsection{What Is Multi-Hop Retrieval and Reasoning?}
\label{subsec:definition}

We define \emph{multi-hop retrieval and reasoning} as a process in which a system acquires and uses information over multiple dependent state transitions, such that at least one later retrieval or reasoning decision is conditioned on information obtained in an earlier step. The definition has three consequences.

First, multi-document processing is not necessarily multi-hop. Retrieving ten passages independently from the original query is still single-stage retrieval. Second, chain-of-thought generation is not necessarily multi-hop retrieval; a model may generate several reasoning sentences without acquiring new information. Third, repeated tool use is not sufficient. The actions must participate in a dependency chain that advances an unresolved information need.

A task may exhibit one or more forms of dependence:
\begin{itemize}[leftmargin=1.35em]
\item \textbf{Entity dependence}: an intermediate entity is required to formulate the next query.
\item \textbf{Relational dependence}: a relation discovered at one step determines the next traversal or source.
\item \textbf{Constraint dependence}: earlier evidence introduces temporal, geographic, numerical, or type constraints for later retrieval.
\item \textbf{Aggregation dependence}: multiple independently found items must be compared, counted, reconciled, or synthesized.
\item \textbf{Verification dependence}: an intermediate claim determines which corroborating or contradictory evidence must be sought.
\end{itemize}

\subsection{Running Examples}
\label{subsec:running_examples}

\paragraph{Bridge question answering.}
The \emph{Kiss and Tell} example in Section~\ref{sec:introduction} requires resolving Shirley Temple before searching for her government position. It illustrates entity dependence and a short, identifiable evidence chain.

\paragraph{Comparative or aggregative questions.}
A question asking which of two organizations reached a milestone first requires locating the relevant date for each organization and then comparing the results. The retrieval branches may be independent initially, but the final decision depends on cross-evidence normalization and comparison.

\paragraph{Open-web research.}
A request to explain why a recent technical product underperformed may require discovering the product version, locating release notes, finding benchmark results, checking issue reports, and reconciling sources with different dates or incentives. The trajectory is longer, evidence quality is heterogeneous, and the output may be a cited report rather than a short answer. Nevertheless, the same functional stages apply.

\subsection{Task Formulation}
\label{subsec:task_formulation}

Let $q$ denote the task, $\mathcal{K}$ the searchable information environment, and $M_i$ the persistent memory or learned state available before task $i$. A within-task trajectory is
\[
\tau_i = (s_0,a_0,o_1,s_1,a_1,o_2,\ldots,s_T,y_i),
\]
where $s_t$ contains the original task, acquired evidence, intermediate reasoning, and execution history; $a_t$ is a search, traversal, reading, verification, or stopping action; $o_{t+1}$ is the resulting observation; and $y_i$ is the final answer or report.

The four functional stages can be written as follows. Initial retrieval obtains entry evidence:
\[
E_1 = \mathcal{R}_0(q,\mathcal{K};M_i).
\]
Follow-up retrieval selects later information needs and obtains additional evidence:
\[
u_t = \pi_{\mathrm{search}}(q,E_{1:t},h_t;M_i), \qquad
E_{t+1}=\mathcal{R}(u_t,\mathcal{K}),
\]
where $h_t$ denotes the trajectory history. Evidence integration constructs a usable representation and output:
\[
C_i = \mathcal{I}(q,E_{1:T},h_T), \qquad
 y_i = \mathcal{G}(q,C_i).
\]
After evaluation signal $f_i$ is observed, feedback learning may update persistent state or policy for future tasks:
\[
M_{i+1}=\mathcal{U}(M_i,\tau_i,f_i),
\qquad
\pi_{i+1}=\mathcal{A}(\pi_i,M_{i+1}).
\]
The formulation separates two timescales. The sequence $E_1,\ldots,E_T$ concerns adaptation within the current task. The update from $M_i$ to $M_{i+1}$ concerns learning across tasks.

A system is multi-hop under this formulation when there exists a step $t>1$ for which the action or evidence relevance depends materially on $E_{1:t-1}$ or $h_{t-1}$, rather than only on the original query $q$. This operational definition can be tested through trajectory inspection, controlled ablations, or counterfactual replacement of intermediate evidence.

\subsection{Taxonomy}
\label{subsec:taxonomy}

The taxonomy follows information flow rather than a specific architecture. Figure~\ref{fig:taxonomy_tree} shows the hierarchy used in Sections~\ref{sec:initial_retrieval}--\ref{sec:feedback_learning}; Table~\ref{tab:taxonomy_overview} summarizes the transformation associated with each category.

\begin{figure*}[tbp]
\centering
\resizebox{0.99\textwidth}{!}{\begin{tikzpicture}[
  x=1cm,y=1cm,
  line cap=round,line join=round,
  branch/.style={draw=black!65,line width=0.55pt},
  titlebox/.style={rounded corners=2pt,draw=taxblue,line width=0.8pt,
                  fill=blue!3,minimum width=15.7cm,minimum height=0.68cm,
                  align=center,font=\bfseries\fontsize{8.2}{9.2}\selectfont},
  category/.style={rounded corners=2pt,draw=taxblue,line width=0.8pt,
                   text width=2.75cm,minimum height=1.05cm,align=center,
                   inner sep=2.5pt,font=\fontsize{7.4}{8.2}\selectfont},
  subcat/.style={rounded corners=2pt,draw=taxblue,line width=0.65pt,
                 text width=3.35cm,minimum height=0.52cm,align=center,
                 inner sep=1.8pt,font=\bfseries\fontsize{6.4}{7.1}\selectfont},
  works/.style={rounded corners=2pt,draw=taxblue,line width=0.55pt,
                text width=6.15cm,minimum height=0.52cm,align=left,
                inner sep=2.0pt,font=\fontsize{5.8}{6.6}\selectfont}
]
\definecolor{taxblue}{RGB}{25,79,120}
\definecolor{initialfill}{RGB}{252,232,232}
\definecolor{followfill}{RGB}{235,237,254}
\definecolor{integrationfill}{RGB}{255,249,218}
\definecolor{feedbackfill}{RGB}{228,250,230}
\definecolor{sectred}{RGB}{205,33,33}

% Left vertical title and common trunk.
\node[titlebox,rotate=90] (root) at (0.42,8.10) {Multi-Hop Retrieval and Reasoning};
\draw[branch] (1.08,0.20) -- (1.08,16.00);

% Main category nodes.
\node[category,fill=initialfill] (c4) at (2.95,14.80)
  {\textbf{Initial Retrieval}\\[-1pt]Section \textcolor{sectred}{4}};
\node[category,fill=followfill] (c5) at (2.95,10.60)
  {\textbf{Follow-up Retrieval}\\[-1pt]Section \textcolor{sectred}{5}};
\node[category,fill=integrationfill] (c6) at (2.95,6.00)
  {\textbf{Evidence Integration}\\[-1pt]Section \textcolor{sectred}{6}};
\node[category,fill=feedbackfill] (c7) at (2.95,1.80)
  {\textbf{Feedback Learning}\\[-1pt]Section \textcolor{sectred}{7}};

\draw[branch] (1.08,14.80) -- (c4.west);
\draw[branch] (1.08,10.60) -- (c5.west);
\draw[branch] (1.08,6.00) -- (c6.west);
\draw[branch] (1.08,1.80) -- (c7.west);

% Section 4: Initial Retrieval.
\node[subcat,fill=initialfill] (s41) at (6.85,16.00)
  {Sparse Retrieval \textcolor{sectred}{(4.1)}};
\node[subcat,fill=initialfill] (s42) at (6.85,15.20)
  {Dense and Hybrid Retrieval \textcolor{sectred}{(4.2)}};
\node[subcat,fill=initialfill] (s43) at (6.85,14.40)
  {Graph Retrieval \textcolor{sectred}{(4.3)}};
\node[subcat,fill=initialfill] (s44) at (6.85,13.60)
  {Reasoning-Aware Retrieval \textcolor{sectred}{(4.4)}};

\node[works,fill=initialfill!45!white] (w41) at (12.35,16.00)
  {BM25~\cite{robertson2009bm25}, COIL~\cite{gao2021coil}, SPLADE~\cite{formal2021splade}.};
\node[works,fill=initialfill!45!white] (w42) at (12.35,15.20)
  {DPR~\cite{karpukhin2020dpr}, ANCE~\cite{xiong2020ance}, Contriever~\cite{izacard2022contriever},\\ColBERT~\cite{khattab2020colbert}, ColBERTv2~\cite{santhanam2022colbertv2}.};
\node[works,fill=initialfill!45!white] (w43) at (12.35,14.40)
  {RAPTOR~\cite{sarthi2024raptor}, GraphRAG~\cite{edge2024graphrag}, HippoRAG~\cite{gutierrez2024hipporag}.};
\node[works,fill=initialfill!45!white] (w44) at (12.35,13.60)
  {HyDE~\cite{gao2023hyde}, REALM~\cite{guu2020realm}.};

\draw[branch] (c4.east) -- (4.80,14.80);
\draw[branch] (4.80,13.60) -- (4.80,16.00);
\foreach \n in {s41,s42,s43,s44}{\draw[branch] (4.80,\n.center|-4.80,0) -- (\n.west);}
\draw[branch] (s41.east)--(w41.west);
\draw[branch] (s42.east)--(w42.west);
\draw[branch] (s43.east)--(w43.west);
\draw[branch] (s44.east)--(w44.west);

% Section 5: Follow-up Retrieval.
\node[subcat,fill=followfill] (s51) at (6.85,12.60)
  {Query Decomposition \textcolor{sectred}{(5.1)}};
\node[subcat,fill=followfill] (s52) at (6.85,11.80)
  {Iterative Retrieval \textcolor{sectred}{(5.2)}};
\node[subcat,fill=followfill] (s53) at (6.85,11.00)
  {Graph and Path Search \textcolor{sectred}{(5.3)}};
\node[subcat,fill=followfill] (s54) at (6.85,10.20)
  {Adaptive Retrieval \textcolor{sectred}{(5.4)}};
\node[subcat,fill=followfill] (s55) at (6.85,9.40)
  {Deep Search \textcolor{sectred}{(5.5)}};
\node[subcat,fill=followfill] (s56) at (6.85,8.60)
  {Feedback-Assisted Retrieval \textcolor{sectred}{(5.6)}};

\node[works,fill=followfill!45!white] (w51) at (12.35,12.60)
  {DecompRC~\cite{min2019decomprc}, BREAK~\cite{wolfson2020break}, Self-Ask~\cite{press2022selfask}.};
\node[works,fill=followfill!45!white] (w52) at (12.35,11.80)
  {MDR~\cite{xiong2021mdr}, Baleen~\cite{khattab2021baleen}, IRCoT~\cite{trivedi2023ircot},\\ITER-RETGEN~\cite{shao2023iterretgen}.};
\node[works,fill=followfill!45!white] (w53) at (12.35,11.00)
  {PullNet~\cite{sun2019pullnet}, HopRetriever~\cite{li2020hopretriever}, ToG~\cite{sun2024tog}.};
\node[works,fill=followfill!45!white] (w54) at (12.35,10.20)
  {FLARE~\cite{jiang2023flare}, Self-RAG~\cite{asai2024selfrag}, Adaptive-RAG~\cite{jeong2024adaptiverag},\\DRAGIN~\cite{su2024dragin}.};
\node[works,fill=followfill!45!white] (w55) at (12.35,9.40)
  {WebGPT~\cite{nakano2021webgpt}, ReAct~\cite{yao2023react}, PaperQA~\cite{lala2023paperqa},\\DeepResearcher~\cite{zheng2025deepresearcher}.};
\node[works,fill=followfill!45!white] (w56) at (12.35,8.60)
  {Search-R1~\cite{jin2025searchr1}, RAG-Gym~\cite{xiong2025raggym}.};

\draw[branch] (c5.east) -- (4.80,10.60);
\draw[branch] (4.80,8.60) -- (4.80,12.60);
\foreach \n in {s51,s52,s53,s54,s55,s56}{\draw[branch] (4.80,\n.center|-4.80,0) -- (\n.west);}
\draw[branch] (s51.east)--(w51.west);
\draw[branch] (s52.east)--(w52.west);
\draw[branch] (s53.east)--(w53.west);
\draw[branch] (s54.east)--(w54.west);
\draw[branch] (s55.east)--(w55.west);
\draw[branch] (s56.east)--(w56.west);

% Section 6: Evidence Integration.
\node[subcat,fill=integrationfill] (s61) at (6.85,7.60)
  {Evidence Selection \textcolor{sectred}{(6.1)}};
\node[subcat,fill=integrationfill] (s62) at (6.85,6.80)
  {Preservation and Compression \textcolor{sectred}{(6.2)}};
\node[subcat,fill=integrationfill] (s63) at (6.85,6.00)
  {Context Construction \textcolor{sectred}{(6.3)}};
\node[subcat,fill=integrationfill] (s64) at (6.85,5.20)
  {Cross-Evidence Reasoning \textcolor{sectred}{(6.4)}};
\node[subcat,fill=integrationfill] (s65) at (6.85,4.40)
  {Long-Form Synthesis \textcolor{sectred}{(6.5)}};

\node[works,fill=integrationfill!45!white] (w61) at (12.35,7.60)
  {SetR~\cite{lee2025setr}, R-CPS~\cite{xin2025readercps}, RankRAG~\cite{yu2024rankrag}.};
\node[works,fill=integrationfill!45!white] (w62) at (12.35,6.80)
  {RECOMP~\cite{xu2023recomp}, LLMLingua~\cite{jiang2023llmlingua},\\LongLLMLingua~\cite{jiang2023longllmlingua}, xRAG~\cite{cheng2024xrag}.};
\node[works,fill=integrationfill!45!white] (w63) at (12.35,6.00)
  {FiD~\cite{izacard2021fid}, Lost in the Middle~\cite{liu2024lostmiddle}, LongRAG~\cite{jiang2024longrag}.};
\node[works,fill=integrationfill!45!white] (w64) at (12.35,5.20)
  {DFGN~\cite{qiu2019dfgn}, HGN~\cite{fang2020hgn}, QA-GNN~\cite{yasunaga2021qagnn},\\GreaseLM~\cite{zhang2022greaselm}.};
\node[works,fill=integrationfill!45!white] (w65) at (12.35,4.40)
  {RARR~\cite{gao2023rarr}, ALCE~\cite{gao2023alce}.};

\draw[branch] (c6.east) -- (4.80,6.00);
\draw[branch] (4.80,4.40) -- (4.80,7.60);
\foreach \n in {s61,s62,s63,s64,s65}{\draw[branch] (4.80,\n.center|-4.80,0) -- (\n.west);}
\draw[branch] (s61.east)--(w61.west);
\draw[branch] (s62.east)--(w62.west);
\draw[branch] (s63.east)--(w63.west);
\draw[branch] (s64.east)--(w64.west);
\draw[branch] (s65.east)--(w65.west);

% Section 7: Feedback Learning.
\node[subcat,fill=feedbackfill] (s71) at (6.85,3.40)
  {Feedback Collection \textcolor{sectred}{(7.1)}};
\node[subcat,fill=feedbackfill] (s72) at (6.85,2.60)
  {Trajectory Diagnosis \textcolor{sectred}{(7.2)}};
\node[subcat,fill=feedbackfill] (s73) at (6.85,1.80)
  {Experience Memory \textcolor{sectred}{(7.3)}};
\node[subcat,fill=feedbackfill] (s74) at (6.85,1.00)
  {Skill Extraction \textcolor{sectred}{(7.4)}};
\node[subcat,fill=feedbackfill] (s75) at (6.85,0.20)
  {Retrieval and Policy Improvement \textcolor{sectred}{(7.5)}};

\node[works,fill=feedbackfill!45!white] (w71) at (12.35,3.40)
  {CRITIC~\cite{gou2024critic}, Self-Refine~\cite{madaan2023selfrefine}.};
\node[works,fill=feedbackfill!45!white] (w72) at (12.35,2.60)
  {Reflexion~\cite{shinn2023reflexion}, ExpeL~\cite{zhao2024expel}.};
\node[works,fill=feedbackfill!45!white] (w73) at (12.35,1.80)
  {Generative Agents~\cite{park2023generativeagents}, LongMem~\cite{wang2023longmem},\\MemGPT~\cite{packer2023memgpt}.};
\node[works,fill=feedbackfill!45!white] (w74) at (12.35,1.00)
  {Voyager~\cite{wang2023voyager}, Memento-Skills~\cite{zhou2026mementoskills}.};
\node[works,fill=feedbackfill!45!white] (w75) at (12.35,0.20)
  {SkillRL~\cite{xia2026skillrl}, Agentic-R~\cite{liu2026agenticr}.};

\draw[branch] (c7.east) -- (4.80,1.80);
\draw[branch] (4.80,0.20) -- (4.80,3.40);
\foreach \n in {s71,s72,s73,s74,s75}{\draw[branch] (4.80,\n.center|-4.80,0) -- (\n.west);}
\draw[branch] (s71.east)--(w71.west);
\draw[branch] (s72.east)--(w72.west);
\draw[branch] (s73.east)--(w73.west);
\draw[branch] (s74.east)--(w74.west);
\draw[branch] (s75.east)--(w75.west);

\end{tikzpicture}
}
\caption{Functional taxonomy of multi-hop retrieval and reasoning. The first three categories describe how a system solves the current task: obtaining entry evidence, acquiring later evidence, and integrating the collected information. Feedback learning operates across tasks by retaining evaluations, trajectories, experiences, skills, or policy updates. The rightmost column lists 60 representative works selected to anchor the taxonomy; the complete \CorpusSize-record corpus and record-level coding are provided in the companion artifact.}
\Description{A taxonomy tree with four branches: Initial Retrieval, Follow-up Retrieval, Evidence Integration, and Feedback Learning. Each branch contains its survey subsections and a rightmost column of representative cited methods.}
\label{fig:taxonomy_tree}
\end{figure*}

\begin{table*}[tbp]
\centering
\caption{Functional taxonomy of multi-hop retrieval and reasoning.}
\label{tab:taxonomy_overview}
\small
\begin{tabularx}{\textwidth}{@{}p{3.0cm}p{3.4cm}Y Y@{}}
\toprule
Category & Transformation & Core question & Main mechanisms \\
\midrule
Initial Retrieval & $q \rightarrow E_1$ & How is a useful entry point obtained from the original task? & Sparse, dense, hybrid, graph, hierarchical, reranking, and reasoning-aware retrieval. \\
Follow-up Retrieval & $(q,E_{1:t}) \rightarrow u_t \rightarrow E_{t+1}$ & What information should be sought next, and when should search stop? & Query decomposition, iterative retrieval, graph/path search, adaptive retrieval, tool use, deep search, and within-task feedback. \\
Evidence Integration & $E_{1:T} \rightarrow C \rightarrow y$ & How are acquired sources selected, preserved, organized, combined, and synthesized? & Evidence selection, compression, context construction, cross-evidence reasoning, grounding, and long-form synthesis. \\
Feedback Learning & $(\tau_i,f_i,M_i) \rightarrow M_{i+1},\pi_{i+1}$ & How does experience from previous tasks improve future behavior? & Feedback collection, trajectory diagnosis, experience memory, skill extraction, retriever learning, and policy optimization. \\
\bottomrule
\end{tabularx}
\end{table*}

\paragraph{Architectural crosswalk.}
The functional categories cut across common system labels. Table~\ref{tab:architectural_crosswalk} shows how typical architectures instantiate the four stages. The entries describe dominant implementations rather than hard membership rules: the same graph, retriever, memory, or language model can serve different functions depending on when it is used and what state it consumes.

\begin{table*}[tbp]
\centering
\caption{Crosswalk from common system families to the four functional categories.}
\label{tab:architectural_crosswalk}
\footnotesize
\begin{tabularx}{\textwidth}{@{}p{2.65cm}Y Y Y Y@{}}
\toprule
System family & Initial retrieval & Follow-up retrieval & Evidence integration & Feedback learning \\
\midrule
Single-shot RAG & One query retrieves a fixed candidate set & Usually absent & Top-$k$ context is passed to a generator & Usually absent \\
Static multi-hop QA & Entry passages or a fixed evidence pool & Fixed or benchmark-defined hops may be executed & Short-answer reader, supporting-fact predictor, or fusion module & Usually absent \\
Iterative or graph retriever & Entry passage, node, neighborhood, or hierarchy & Learned hop expansion, path search, or graph traversal & Reader or graph reasoning over the acquired path/subgraph & Usually absent unless trajectories are retained \\
Adaptive RAG or search agent & Initial query, tool choice, or source selection & Dynamic retrieve/read/verify/stop actions & Context construction and answer synthesis from the current trajectory & Within-task feedback by default; persistent learning is optional \\
Deep search or deep research & Broad source discovery and task planning & Long-horizon open-web search with variable branching and stopping & Multi-source report generation, citation, and contradiction handling & Optional experience, memory, or policy updates across research tasks \\
Experience- or skill-enhanced system & Entry behavior may be conditioned on retrieved memories or skills & Later actions reuse lessons, procedures, or a learned policy & Skills may prescribe evidence processing and synthesis & Defining persistent update through reflection, memory, skills, retriever training, or policy learning \\
\bottomrule
\end{tabularx}
\end{table*}

The taxonomy assigns each paper a primary category according to the function most directly targeted by its main contribution, while secondary tags record cross-category effects. For example, Self-RAG performs adaptive retrieval and critique within a task and is therefore primarily discussed under follow-up retrieval, even though its critique mechanism also affects generation. A skill library extracted from completed search trajectories is primarily feedback learning because its defining contribution is reuse across future tasks. This primary-label rule prevents the corpus from double-counting methods while preserving architectural crosswalks.

The categories also define chapter boundaries. Pure query--document reranking belongs to initial retrieval. A later query generated from retrieved evidence belongs to follow-up retrieval. Selection or ordering of an already acquired evidence set belongs to evidence integration. Memory or reflection belongs to feedback learning only when it changes future task behavior; transient scratchpads and current-task observations remain part of the within-task trajectory.

\section{Initial Retrieval}
\label{sec:initial_retrieval}

Initial retrieval establishes the entry evidence available before any task-specific intermediate result has been resolved. Its input is the original query or task description, and its output is a candidate set, ranked list, graph neighborhood, or structured entry point for later search and reasoning. In multi-hop settings, a strong initial retriever must balance local relevance with the probability of exposing bridge information that enables later hops.

\subsection{Sparse Retrieval}
\label{subsec:sparse_retrieval}
% Cover BM25, learned sparse retrieval, query expansion, lexical robustness,
% and the role of sparse signals in hybrid multi-hop systems.

\subsection{Dense and Hybrid Retrieval}
\label{subsec:dense_hybrid_retrieval}
% Cover DPR, ANCE, Contriever, ColBERT, E5/BGE-style embeddings,
% hybrid retrieval, reranking, and reader-aware retrieval refinement.

\subsection{Graph Retrieval}
\label{subsec:graph_retrieval}
% Cover graph construction, subgraph retrieval, hierarchical retrieval,
% GraphRAG/RAPTOR/HippoRAG-style entry points, and structured evidence access.

\subsection{Reasoning-Aware Retrieval}
\label{subsec:reasoning_aware_retrieval}
% Cover retrievers trained or prompted with reasoning signals, synthetic
% explanations, hypothetical documents, process rewards, or downstream reader utility.

\paragraph{Chapter boundary.}
Methods remain in this chapter when they operate primarily on the original task or a fixed reformulation of it. Once retrieved evidence is used to formulate the next information need, the method is discussed in Section~\ref{sec:follow_up_retrieval}.

\section{Follow-up Retrieval}
\label{sec:follow_up_retrieval}

Follow-up retrieval begins when the next information need depends on evidence, entities, constraints, or uncertainty produced earlier in the trajectory. This chapter therefore focuses on search-state transitions rather than first-stage ranking alone. The key design questions are how to represent the current state, how to generate or select the next query or action, how to control branching and stopping, and how to recover from earlier retrieval errors.

\subsection{Query Decomposition}
\label{subsec:query_decomposition}
% Cover explicit subquestions, decomposition supervision, modular prompting,
% decomposition faithfulness, and decomposition error propagation.

\subsection{Iterative Retrieval}
\label{subsec:iterative_retrieval}
% Cover GoldEn Retriever, MDR, Baleen, IRCoT, iterative retrieval-generation,
% and hop-conditioned dense retrieval.

\subsection{Graph and Path Search}
\label{subsec:graph_path_search}
% Cover iterative graph expansion, evidence-path retrieval, beam search,
% KG traversal, and graph-guided LLM search.

\subsection{Adaptive Retrieval}
\label{subsec:adaptive_retrieval}
% Cover uncertainty-triggered retrieval, complexity routing, retrieve/read/stop
% policies, self-reflection within a task, and corrective retrieval.

\subsection{Deep Search}
\label{subsec:deep_search}
% Cover browser and open-web search agents, longer-horizon research trajectories,
% source discovery, stopping, and report-oriented information acquisition.

\subsection{Feedback-Assisted Retrieval}
\label{subsec:feedback_assisted_retrieval}
% Cover within-task process feedback, reward-guided search, verifier-guided
% retrieval, and search policies trained from trajectory-level signals.

\paragraph{Within-task versus cross-task feedback.}
Feedback belongs in this chapter when it changes a later action in the same task. Feedback that is retained to improve subsequent tasks is treated as feedback learning in Section~\ref{sec:feedback_learning}.

\input{sections/06_evidence_organization_reasoning}
\section{Feedback Learning}
\label{sec:feedback_learning}

Feedback learning studies how completed trajectories change behavior on future tasks. It differs from adaptive retrieval within a single task: the defining update persists beyond the current episode. The retained object may be a reflection, an experience record, a memory item, a reusable procedure, a skill, retriever supervision, or an updated search policy.

\subsection{Feedback Collection}
\label{subsec:feedback_collection}
% Cover outcome rewards, process rewards, verifier signals, human feedback,
% tool feedback, environment feedback, and automatic critics.

\subsection{Trajectory Diagnosis}
\label{subsec:trajectory_diagnosis}
% Cover failure attribution, reflection, comparison of successful and failed
% trajectories, credit assignment, and diagnosis granularity.

\subsection{Experience Memory}
\label{subsec:experience_memory}
% Cover episodic experience stores, distilled lessons, retrieval of prior
% trajectories, memory consolidation, and interference or staleness.

\subsection{Skill Extraction}
\label{subsec:skill_extraction}
% Cover executable skills, textual procedures, skill libraries, automatic
% skill revision, skill routing, and compositional reuse.

\subsection{Retrieval and Policy Improvement}
\label{subsec:policy_improvement}
% Cover reinforcement-learned search policies, process-supervised retrievers,
% retriever updates from downstream trajectories, and continual adaptation.

\paragraph{Inclusion rule.}
General memory or skill systems are included only when their retained experience directly improves retrieval-grounded multi-hop behavior. A long-term memory module that merely stores user facts, or an embodied skill unrelated to information acquisition, is outside the primary scope.

\section{Evaluation}
\label{sec:evaluation}

Evaluation should match the component or system claim being made. A gain in final-answer accuracy may result from better retrieval, a stronger reader, additional context, memorized knowledge, or a different stopping policy. Multi-hop evaluation therefore requires both end-to-end outcomes and process-level measurements that expose where information was found, lost, ignored, or reused. The four functional categories in this survey provide a claim-aligned evaluation template: initial retrieval concerns the entry pool, follow-up retrieval concerns trajectories, evidence integration concerns preservation and use, and feedback learning concerns persistent improvement across tasks.

\subsection{Retrieval Evaluation}
\label{subsec:retrieval_evaluation}

Standard retrieval metrics include Recall@$K$, precision, mean reciprocal rank, and nDCG. For multi-hop tasks, these metrics should be computed over different evidence roles. \emph{Bridge recall} measures whether evidence needed to formulate a later query is available. \emph{Answer-evidence recall} measures whether answer-bearing sources are available. \emph{Complete-chain recall} measures whether all required units are present simultaneously. A system can have high document recall while failing complete-chain recall if it repeatedly retrieves the same hop.

Initial and follow-up retrieval should also be evaluated over budget curves. Reporting a single $K$ hides whether a method improves retrieval or simply requires a larger candidate pool. Useful curves vary candidate depth, number of retrieval rounds, total retrieved tokens, and unique sources. For graph systems, analogous quantities include subgraph size, path coverage, and expansion depth. When the dataset supplies supporting facts, as in HotpotQA and MuSiQue~\cite{yang2018hotpotqa,trivedi2022musique}, retrieval metrics should be computed before reader filtering so that retrieval and evidence use remain separable.

Candidate-pool control is essential for reranking and selection claims. Comparisons should fix the first-stage pool, context budget, and reader. Oracle analyses can estimate the headroom available if the correct evidence is selected, but oracle evidence should not be mixed with open-domain results. Datasets may also contain disconnected reasoning shortcuts; answer accuracy should therefore be interpreted together with evidence coverage and shortcut-resistant subsets~\cite{trivedi2020nohotpot}.

\subsection{Search Trajectory Evaluation}
\label{subsec:trajectory_evaluation}

Follow-up retrieval produces a sequence rather than a single ranking. Trajectory evaluation should measure whether each action advances an unresolved information need. For decomposition methods, relevant metrics include subquestion validity, dependency preservation, answerability, and coverage of the original task. For iterative retrieval, useful metrics include hop success, path recall, recovery after an incorrect hop, branching factor, and the fraction of trajectories that terminate with a complete support set.

Adaptive and agentic systems require action-level diagnostics. These include retrieval-trigger precision and recall, invalid-action rate, query duplication, source revisit rate, premature stopping, unnecessary retrieval, and the number of tool or browser actions. Cost should be reported in at least two units: environment interactions, such as searches and pages opened, and model consumption, such as tokens or inference calls. Wall-clock time and monetary API cost are useful but hardware- and provider-dependent.

Trajectory quality cannot be reduced to length. A shorter trajectory may be efficient or may stop before the task is solved; a longer one may be wasteful or may correctly verify conflicting evidence. Cost-adjusted metrics should therefore condition on task success, for example successful tasks per search call or answer quality under a fixed action budget. Releasing trajectories, query logs, and stopping decisions enables later analysis and should be considered part of the evaluation artifact when privacy and licensing permit.

\subsection{Evidence Integration Evaluation}
\label{subsec:integration_evaluation}

Evidence-integration evaluation asks whether acquired information survives selection and compression, is presented in a usable form, and causally supports the output. Membership and ordering must be separated. To test a selector, fix the candidate pool and presentation policy. To test ordering, fix the evidence set. To test compression, compare the compressed representation with the same uncompressed sources under matched reader conditions.

Selection metrics include complete-support preservation, redundancy, diversity, contradiction rate, and reader utility. Compression metrics include compression ratio, token and latency reduction, support retention, relation retention, and provenance retention. Context-construction studies should sweep context length and evidence position because long-context models can degrade as additional retrieved passages are added~\cite{liu2024lostmiddle,jin2025longcontextrag}.

Grounding requires claim-level checks. Citation correctness asks whether a cited source supports the associated claim; citation completeness asks whether claims that require evidence are cited. ALCE formalizes several aspects of citation-grounded generation~\cite{gao2023alce}. Supporting-fact or rationale overlap is useful when annotations exist, but generated rationales may be post hoc. Stronger tests intervene on evidence: remove a required source, replace a relation, permute dependent evidence, or introduce a contradictory source. If the output does not respond appropriately, the system may not be using the claimed evidence.

Long-form synthesis adds coverage, coherence, temporal consistency, and source-quality dimensions. Overall report scores should be accompanied by claim--source mappings and per-claim judgments. Human evaluators should have access to the cited sources; otherwise, fluency and topical plausibility can dominate the score.

\subsection{Feedback Learning Evaluation}
\label{subsec:feedback_evaluation}

Feedback-learning claims require a sequence of tasks or repeated rounds with persistent state. A successful retry on the same task is insufficient evidence of cross-task learning unless the retained update is isolated. Evaluation should distinguish \emph{within-task recovery}, \emph{same-distribution transfer}, and \emph{out-of-distribution transfer}.

For experience memories, report memory write rate, retrieval hit rate, retrieved-memory relevance, context overhead, and performance with memory retrieval disabled. For skill libraries, report skill selection accuracy, execution success, reuse frequency, library growth, routing cost, and composition failures. For learned retrievers or policies, report learning curves, held-out-task performance, reward definitions, and ablations that freeze each jointly evolving component.

Negative transfer and forgetting are first-class outcomes. A growing memory can retrieve stale or conflicting lessons; a revised skill can improve one task family while damaging another; and a retriever trained on agent-generated queries can overfit to the current policy. Evaluation should include chronological tests, memory-size sweeps, deletion or replacement interventions, and performance on tasks not represented in the retained experience. The strongest evidence links a specific retained object to a changed action and a downstream improvement on a later task.

\subsection{Benchmarks}
\label{subsec:benchmarks}

Benchmarks differ in which parts of the process they make observable. HotpotQA supplies supporting facts and diverse bridge and comparison questions~\cite{yang2018hotpotqa}. MuSiQue constructs questions from composed single-hop problems and reduces several disconnected shortcuts~\cite{trivedi2022musique}. HybridQA and MultiModalQA expand the evidence substrate to tables, text, and images~\cite{chen2020hybridqa,talmor2021multimodalqa}. These datasets are useful for controlled component evaluation because evidence sets or intermediate structure are partly annotated.

Open-web and deep-search evaluations trade annotation density for ecological validity. Browser-assisted systems such as WebGPT operate over changing web content and action interfaces~\cite{nakano2021webgpt}. Such benchmarks should snapshot sources or log timestamps, search-provider settings, and page content so that results can be reproduced. Long-form and citation-oriented benchmarks should separate retrieval, source use, and writing quality rather than collapsing them into one judge score.

No benchmark covers all four categories. A task with gold supporting facts may reveal retrieval and integration errors but provide no repeated-task signal for feedback learning. An agent benchmark may expose trajectories but lack evidence annotations. A continual-learning benchmark may measure transfer without source grounding. Table~\ref{tab:evaluation_alignment} therefore presents a minimum alignment between claims, metrics, and controls.

\begin{table*}[tbp]
\centering
\caption{Claim-aligned evaluation for multi-hop retrieval and reasoning.}
\label{tab:evaluation_alignment}
\small
\begin{tabularx}{\textwidth}{@{}p{2.8cm}p{4.1cm}Y Y@{}}
\toprule
Claim & Minimum direct measurements & Required controls & Common confound \\
\midrule
Better initial retrieval & Bridge recall, answer-evidence recall, complete-chain Recall@$K$ & Fixed corpus, candidate depth, and query & Stronger reader or larger candidate pool \\
Better follow-up retrieval & Hop/path success, query/action trace, stopping, total search cost & Fixed entry pool and action budget & More retrieval calls or hidden tool access \\
Better evidence integration & Support preservation, ordering/compression test, grounding or citation metrics & Fixed evidence membership, reader, and context budget & Selection and reader changes occur together \\
Better long-form synthesis & Claim coverage, citation correctness/completeness, contradiction handling & Source snapshot and claim-level evaluation & Fluency dominates judge score \\
Better feedback learning & Learning curve, memory/skill use, held-out transfer, negative transfer & Persistent-state ablation and matched context cost & Repeated sampling or within-task retry \\
Better end-to-end system & Answer/report quality plus retrieval, trajectory, grounding, and cost & Same data access, tools, model budget, and stopping limits & Unreported differences in resources or model scale \\
\bottomrule
\end{tabularx}
\end{table*}

\paragraph{Diagnostic audit.}
The companion audit uses a stratified \AuditSize-paper selection across the four functional categories. Each record is coded only after a full-text source passage, table, figure, or intervention is located. The codebook distinguishes \emph{reported}, \emph{not reported}, \emph{not applicable}, and \emph{unclear}; absence is never inferred from a search snippet or abstract. This design allows the survey to summarize evaluation practice without treating the larger discovery corpus as fully audited evidence.

\paragraph{Summary.}
Multi-hop evaluation is credible when it reveals both outcome and process. The required measurements depend on the claim, but the general rule is stable: fix neighboring components, report budgets and trajectories, and test whether the evidence or experience attributed to success is actually used.

\section{Future Directions}
\label{sec:future}

The functional taxonomy exposes several gaps that are difficult to see from architecture-level comparisons alone. The central problem is no longer whether a system can issue several searches, but whether it can maintain a useful information state over long horizons, preserve the dependencies required by later reasoning, and improve without accumulating brittle or stale behavior.

\subsection{Retrieval for Future Utility}
Most retrievers are trained to estimate immediate query--document relevance. Multi-hop systems additionally need evidence whose value is conditional: a bridge passage may not answer the current question, yet may identify the entity, relation, constraint, or vocabulary required for a later search. Future work should therefore distinguish \emph{answer evidence}, \emph{bridge evidence}, and \emph{search-enabling evidence}, and learn objectives that value their different roles.

A key research question is how to estimate future utility without requiring a fully annotated gold chain. Possible signals include successful downstream queries, complete-chain recovery, counterfactual removal of bridge passages, and process rewards tied to subsequent search progress. Evaluation should report whether the method improves entry evidence, later-hop reachability, or both.

\subsection{State Representation and Error Recovery in Long-Horizon Search}
A long trajectory cannot be controlled reliably from an unstructured transcript alone. The system must represent what has been established, what remains uncertain, which sources support each claim, which queries have already failed, and which branches remain open. State representations may combine textual notes, entity--relation structures, query histories, uncertainty estimates, and source provenance.

Long-horizon systems also require recovery mechanisms. One incorrect entity resolution can redirect all later searches, while repeated reformulation can amplify a false assumption. Future methods should detect inconsistent intermediate states, revisit earlier decisions, compare alternative branches, and stop when additional search is unlikely to change the result. These capabilities should be evaluated through controlled perturbations of intermediate evidence rather than only through naturally occurring failures.

\subsection{Joint Optimization without Losing Component Accountability}
Modern systems increasingly optimize retrieval, reasoning, tool use, and generation jointly. End-to-end optimization can improve final task reward, but it can also obscure which component changed and whether the system learned a shortcut. A useful research direction is modular or constrained joint learning in which search, evidence integration, and generation share feedback while retaining stage-level observability.

This requires reporting interfaces between components: the candidate pool before and after search control, the evidence set before and after integration, and the trajectory before and after feedback updates. Without these controls, a gain attributed to retrieval may instead result from a stronger reader, a larger context, or additional model calls.

\subsection{Evidence Preservation with Explicit Provenance}
Compression and note taking should preserve not only semantic content but also the dependency structure needed by later reasoning. Important details include entity identity, negation, numerical values, temporal qualifiers, uncertainty, and the source from which each statement was derived. A compact note that removes provenance or merges conflicting claims may be efficient but unusable for grounded synthesis.

Future systems should represent compressed evidence as structured claims linked to source spans. Compression evaluation should combine token reduction with relation preservation, answerability under controlled evidence, and citation recovery. This would make it possible to separate harmless wording compression from deletion of a reasoning-critical bridge.

\subsection{Faithful Long-Form Synthesis}
Deep research shifts the output from a short answer to a report containing many claims with different evidential status. The system must decide which claims are directly stated by sources, which are derived by combining sources, and which remain speculative. It must also reconcile contradictions, avoid citing a source for a stronger claim than it supports, and preserve temporal context.

A promising direction is claim--source graph construction before report generation. Such a representation could support coverage checks, contradiction detection, citation placement, and targeted re-search for unsupported sections. Evaluation should measure claim-level support and completeness in addition to document-level citation counts or holistic report ratings.

\subsection{Reusable Experience and Skills}
Feedback learning introduces a second horizon: the system must improve across tasks rather than only within one trajectory. Current experience stores and skill libraries often lack clear policies for retention, retrieval, revision, and deletion. A successful trajectory may encode a domain-specific shortcut that transfers poorly, while an old skill may remain highly ranked after the environment changes.

Future work should evaluate experience reuse on genuinely new tasks, report the retrieval accuracy of memories or skills, and measure negative transfer and forgetting. Skill representations should expose their preconditions, expected effects, evidence requirements, and failure history. Versioning and conflict resolution are necessary if independently learned procedures are to be composed safely.

\subsection{Benchmarks for Open and Repeated Information Seeking}
Most multi-hop benchmarks evaluate a fixed corpus and a single episode. Open-web research introduces changing documents, heterogeneous source quality, access cost, and incomplete gold evidence. Feedback learning additionally requires repeated tasks and a temporal split between experience acquisition and later evaluation.

New benchmarks should release or reconstruct action traces, source visits, retrieved evidence, stopping decisions, and costs. They should distinguish static knowledge access from live-web search, and current-task adaptation from improvements caused by prior experience. For long-form outputs, benchmark annotations should connect claims to source spans rather than provide only a global score.

\subsection{Safety and Governance of Autonomous Information Seeking}
Adaptive retrieval expands the attack surface of a language model. Malicious pages can manipulate queries or tool calls; biased source selection can systematically omit perspectives; persistent memory can retain private or poisoned information; and long-running agents can incur unbounded cost. These risks arise during acquisition and learning, not only in the final text.

Safety evaluation should therefore include source trust, prompt-injection resistance, privacy controls, memory deletion, budget enforcement, and auditability of search actions. Persistent feedback mechanisms require particular scrutiny because a single poisoned trajectory can affect many future tasks. The same process-level logging proposed for scientific evaluation is also necessary for governance and incident analysis.

\section{Conclusion}
\label{sec:conclusion}

Multi-hop reasoning is fundamentally an information-acquisition problem rather than merely a generation problem. A system must determine what is already known, identify what remains unresolved, decide what information should be sought next, and combine the resulting evidence into a supported answer or report. Longer outputs, larger context windows, or more tool calls do not by themselves constitute multi-hop capability. The defining property is dependency: later retrieval and reasoning steps become meaningful because of information acquired earlier.

This survey organized the field into four functions. \emph{Initial retrieval} establishes an entry point from the original query. \emph{Follow-up retrieval} converts intermediate evidence into new information needs and search actions. \emph{Evidence integration} preserves, organizes, and combines information gathered across hops. \emph{Feedback learning} uses evaluations and prior trajectories to improve future retrieval and reasoning. The first three functions explain how a system completes the current task; the fourth explains how repeated task execution can produce durable improvement. This separation makes it possible to compare multi-hop QA pipelines, graph retrieval, adaptive RAG, agentic search, deep research, memory, and skill-based systems without treating them as unrelated research areas or reducing them to a single architectural template.

The evaluation implications are equally important. Final-answer accuracy cannot reveal whether the system found the required evidence, followed a valid search path, preserved key relations under a context budget, actually used the cited sources, stopped efficiently, or learned a reusable strategy. Progress therefore requires process-level evaluation of retrieval depth, hop completion, trajectory quality, evidence membership and ordering, grounding, cost, and cross-task improvement. Benchmarks and reporting practices should make these intermediate objects observable whenever a paper claims to improve them.

The next generation of multi-hop systems will be better characterized as adaptive information-seeking systems than as a retriever attached to a language model. Such systems will need to search, reason, verify, manage context, preserve provenance, and reuse experience under changing information environments. The central research objective is not simply to make models produce more elaborate reasoning traces, but to make them reliably acquire the right information, recognize unresolved uncertainty, integrate evidence without losing support, and improve their behavior from previous attempts.


\bibliographystyle{ACM-Reference-Format}
\bibliography{references}

\end{document}
